% ---
% Subsubsección 3.5.6 (nueva, 2026-10-18). Experimento X6 (D-038, H-047, H-048).
% La identidad de la ec. eq:kpq_pauli se comprobó numéricamente con Qiskit (estados aleatorios de 3
% qubits, partial_trace frente a valores esperados de Pauli): diferencia de 1e-16.
% Implementación en X6: valores esperados <X_k>, <Y_k>, <Z_k> de los 12 qubits del ZZFeatureMap con
% c = 1 y kernel gaussiano sobre ellos, con cinco anchos de banda.
% Cada atribución lleva su respaldo verificable con Code/tools/verificar_citas_fichas.py.
% ---

\subsubsection{Kernel cuántico proyectado}
\label{sec:kernel_proyectado}

% respaldo: «To circumvent this setback, we propose a family of projected quantum kernels as a solution. These kernels work by projecting the quantum states to an approximate classical representation, e.g., using reduced physical observables or classical shadows25,27,42–44.» (huang2021power, p. 6)
El kernel de fidelidad compara estados en el espacio de Hilbert completo, de dimensión $2^n$, donde, como se vio en la sección~\ref{sec:limites_ventaja}, las entradas tienden a quedar demasiado lejos entre sí. Huang et al. propusieron los kernels cuánticos proyectados, que primero proyectan cada estado a una representación clásica aproximada, por ejemplo mediante observables reducidos, y después comparan esas representaciones \cite{huang2021power}.

% respaldo: «One of the simplest forms of projected quantum kernel is to measure the one-particle reduced density matrix (1-RDM) on all qubits for the encoded state» (huang2021power, p. 6)
La forma más simple usa la matriz de densidad reducida de cada qubit \cite{huang2021power}. Si $\rho(\mathbf{x}) = |\phi(\mathbf{x})\rangle\langle\phi(\mathbf{x})|$ es el estado codificado y $\rho_k(\mathbf{x}) = \operatorname{Tr}_{j \neq k}[\rho(\mathbf{x})]$ el del qubit $k$, que se obtiene trazando los demás, el kernel es
\begin{equation}
  k^{\mathrm{PQ}}(\mathbf{x}, \mathbf{x}') = \exp\!\left(-\gamma \sum_{k=1}^{n} \bigl\lVert \rho_k(\mathbf{x}) - \rho_k(\mathbf{x}') \bigr\rVert_F^2\right).
  \label{eq:kpq}
\end{equation}

Esta expresión tiene una lectura directa en términos de mediciones. El estado de un qubit se escribe como $\rho_k = \tfrac{1}{2}\bigl(I + \sum_{P} \langle P_k \rangle\, P\bigr)$, con $P \in \{X, Y, Z\}$, y como $\operatorname{Tr}(PQ) = 2\,\delta_{PQ}$ para las matrices de Pauli,
\begin{equation}
  \bigl\lVert \rho_k(\mathbf{x}) - \rho_k(\mathbf{x}') \bigr\rVert_F^2 = \frac{1}{2} \sum_{P \in \{X,Y,Z\}} \bigl(\langle P_k \rangle_{\mathbf{x}} - \langle P_k \rangle_{\mathbf{x}'}\bigr)^2 .
  \label{eq:kpq_pauli}
\end{equation}
El kernel proyectado es, por tanto, un kernel gaussiano con parámetro $\gamma/2$ sobre el vector de $3n$ valores esperados $\langle X_k \rangle$, $\langle Y_k \rangle$, $\langle Z_k \rangle$ de cada qubit. Se parece al \textit{quantum embedding} de la sección~\ref{sec:quantum_embedding}, con dos diferencias: aquí se leen los tres ejes de cada qubit, y se leen directamente sobre el estado del \textit{feature map}, sin \textit{ansatz}. Su parámetro $\gamma$ es un ancho de banda que se elige como en un kernel clásico.

% Añadido el 2026-10-23 (D-043, D-044): resume las tres formas de usar el estado del feature map que
% aparecen en el reporte. Las cifras de cada una están en la tabla tab:tres_lecturas de la sección 6.4.
Con este kernel, el reporte usa el estado de un mismo \textit{feature map} de tres maneras, que
difieren en qué se obtiene del estado y en cómo se comparan dos lesiones (tabla~\ref{tab:lecturas_estado}).
La distinción importa porque las tres pueden llevar a conclusiones distintas sobre la misma codificación.

\begin{table}[htbp]
\centering
\small
\caption{Tres formas de leer el estado $|\phi(\mathbf{x})\rangle$ de un mismo \textit{feature map}.}
\label{tab:lecturas_estado}
\begin{adjustbox}{max width=\linewidth}
\begin{tabular}{p{0.2\linewidth} p{0.25\linewidth} p{0.25\linewidth} p{0.25\linewidth}}
\toprule
 & \textbf{Valores esperados $\langle Z_i\rangle$} & \textbf{Kernel de fidelidad} & \textbf{Kernel proyectado} \\
\midrule
Circuito & \textit{Feature map} seguido del \textit{ansatz} $V(\boldsymbol\theta)$ & \textit{Feature map} & \textit{Feature map} \\
Qué se obtiene & $\langle Z_i\rangle$ de cada qubit tras el \textit{ansatz} & El solapamiento $|\langle\phi(\mathbf{x})|\phi(\mathbf{x}')\rangle|^2$ entre los estados completos & $\langle X_k\rangle$, $\langle Y_k\rangle$ y $\langle Z_k\rangle$ de cada qubit \\
Resultado & Un vector de $n$ valores por lesión & Un valor por par de lesiones & Un vector de $3n$ valores por lesión, comparado con un kernel gaussiano \\
Escala de la comparación & Local: un eje por qubit & Global: las $2^n$ amplitudes & Local: los tres ejes de cada qubit \\
Depende de $\boldsymbol\theta$ & Sí & No: $V(\boldsymbol\theta)$ se cancela & No: no hay \textit{ansatz} \\
Clasificador & Perceptrón multicapa & SVM (QSVM) & SVM \\
Medición en un procesador & Una base ($Z$) para todos los qubits & Un circuito \textit{compute-uncompute} por par de lesiones & Tres bases ($X$, $Y$, $Z$) \\
Dónde se usa & C4 y C5 (M5 y M7) & M6, X4 y X6 & X6 y X11 \\
\bottomrule
\end{tabular}
\end{adjustbox}
\end{table}

% respaldo: «In numerical experiments, we find that the classical projection increases rather than decreases the geometric difference with classical ML models.» (huang2021power, p. 6)
En los experimentos numéricos de Huang et al., esta proyección aumentó, en lugar de reducir, la \textit{geometric difference} con los modelos clásicos \cite{huang2021power}. En este trabajo se evalúa en un experimento complementario, sobre los estados del ZZFeatureMap con el factor de escala del diseño (sección~\ref{sec:qsvm_ventaja}).
